\chapter{Urinary Frequency}

\section*{Definition}
Increased number voids day without necessarily increasing total urine volume.

\section*{Pathophysiology}
Lower urinary tract symptoms benign prostatic hyperplasia UTIs diabetes bladder overactivity common causes.

\section*{Etiology}
\begin{itemize}
  \item UTI
  \item Diabetes mellitus
  \item Benign prostatic hyperplasia
  \item Overactive bladder
  \item Fluid intake increase
\end{itemize}

\section*{Indications for Evaluation}
Seek care if:
\begin{itemize}
  \item Severe or worsening symptoms
  \item Symptoms lasting more than Evaluated if new bothersome
  \item Accompanied by other concerning signs
\end{itemize}

\section*{Related Symptoms}
\begin{itemize}
  \item Urgency nocturia dysuria hematuria pain burning sensation
\end{itemize}
\textbf{Note:} Always consider serious underlying conditions when evaluating persistent urinary frequency.

\section*{Self-Care / Home Management}
\begin{itemize}
  \item Drink adequate water throughout the day unless fluid restriction is medically indicated.
  \item Avoid bladder irritants including caffeine, alcohol, acidic foods, and artificial sweeteners.
  \item Practice good hygiene and urinate promptly when the urge arises to prevent urinary stasis.
  \item Apply a heating pad to the lower abdomen or back for comfort during urinary discomfort.
  \item Seek medical attention for fever, flank pain, visible blood in urine, or difficulty urinating.
\end{itemize}

\section*{Diagnostic Approach}
\begin{itemize}
  \item History covers urinary symptoms (frequency, urgency, dysuria, hesitancy, stream changes), fever, and flank pain.
  \item Urinalysis with microscopy and culture is the first-line diagnostic test for urinary symptoms.
  \item Imaging (renal ultrasound, CT urogram) evaluates for stones, obstruction, structural abnormalities, or masses.
  \item Urodynamic studies or cystoscopy may be indicated for complex or refractory voiding dysfunction.
\end{itemize}

\section*{Treatment / Management Overview}
\begin{itemize}
  \item Urinary tract infections are treated with targeted antibiotics based on culture sensitivity results.
  \item Nephrolithiasis management includes hydration, analgesics, medical expulsive therapy, or urologic intervention for large stones.
  \item Voiding dysfunction may respond to behavioral modification, pelvic floor therapy, or pharmacotherapy.
  \item Urology referral is indicated for recurrent infections, hematuria, obstructive symptoms, or suspected malignancy.
\end{itemize}

\clearpage
